% !TeX root = ../main.tex
\section{Hạn Chế}
\label{sec:limitations}

Phần mềm hiện chưa có chế độ online nhiều người chơi từ xa. Hàm đánh
giá của AI chỉ dùng hai yếu tố (quân số và mobility); các yếu tố chiến
lược sâu hơn như kiểm soát trung tâm, đe dọa nhiều hướng, hay giá trị
vị trí (positional value) chưa được mô hình hóa.

Mô hình toán học hình thức ở Phụ lục~\ref{app:formal} chỉ đầy đủ cho chế
độ 2 người. Phần mở rộng sang 3 và 4 người chỉ được đặc tả ở mức cấu hình
(kích thước bàn, vị trí xuất phát---Phụ lục~\ref{app:notation}) chứ chưa
được phân tích lý thuyết, đặc biệt là các tính chất cân bằng của ván nhiều
người. Đáng chú ý, cấu hình xuất phát hiện tại của chế độ 3 người
\emph{không cân đối}: hai phe có 13 quân còn phe thứ ba có 14 quân---một
bất đối xứng cần được khảo sát và có thể hiệu chỉnh trong nghiên cứu cân
bằng tương lai.

Giải đấu bot-vs-bot ở Mục~\ref{sec:eval} mới chỉ so các độ sâu tìm kiếm
trong chế độ 2 người. Phép đo lợi thế tay đi đầu
(Mục~\ref{sec:eval_tournament}) cũng chỉ ở chế độ 2 người và tại một độ
sâu cố định ($d{=}2$, nơi bên đi sau tỏ ra có lợi thế); chưa có nghiên cứu
định lượng về độ cân bằng giữa các màu trong chế độ 3 và 4 người---câu hỏi
liệu màu đi trước/đi sau có lợi thế hệ thống ở các độ sâu khác, hay chiến
lược liên minh ngầm có ổn định, còn để mở.

Về điều kiện kết thúc, chế độ 2 người \emph{không} có luật hoà do lặp thế:
trên lý thuyết hai phe có thể chơi vô hạn mà không phe nào bị phong toả
(cài đặt đặt \texttt{max\_stalemate} cho riêng chế độ 3--4 người). Cơ chế
phạt lặp thế của AI (Mục~\ref{ssec:pitfalls}) làm điều này hiếm trong thực
tế nhưng không loại trừ về mặt luật; một luật hoà do lặp thế là hướng bổ
sung tự nhiên.
